\section{Abel rigidity and the flat equal-edge Gaussian comparison}
\label{sec:abel-gaussian-20261016}
\noindent\textbf{Scope.} Claims \texttt{BGS-TRACE-092} and
\texttt{BGS-TRACE-093}. The canonical full proofs are
\path{mathematics/abel_clock_rigidity_20261016.md} and
\path{mathematics/gaussian_equal_edge_20261016.md}. This section records
their statements and proof mechanisms. They do not establish the signed
Heisenberg hierarchy of the maintained fixed propagator. The Gaussian
benchmark and its double-scaling function are classical
\cite{KPSS2002,ByunForrester2023}.

\subsection{A positive decomposition with a geometric remainder}
For eigenvalues $e^{2\pi i\theta_j}$ of $U\in U(N)$ put
$x_j=N\theta_j$, $T_m=\operatorname{Tr}U^m$ and
\[
 L_N(s)=N^{-2}\sum_{m\ge1}e^{-2sm/N}|T_m|^2,\qquad
 \delta_N(s)=L_N(s)-(e^{2s}-1)^{-1}.
\]
For $s>0$ define
\[
 f_s(x)=\frac{s}{s^2+\pi^2x^2},\qquad
 q_s(x)=\frac{s^3}{\sinh^2s}
      \frac{(\sin(\pi x)/(\pi x))^2}{s^2+\pi^2x^2},
 \qquad H_s=f_s-q_s.
\]
Continuous values are used at zero. With Fourier convention
$\widehat g(t)=\int g(x)e^{-2\pi itx}\,dx$,
\[
 \widehat H_s(t)=
 \begin{cases}
 [\sinh(2s(1-|t|))-2s(1-|t|)]/(2\sinh^2s),&|t|\le1,\\
 0,&|t|>1.
 \end{cases}
\]
\begin{theorem}[Exact positive Abel decomposition]
For every $N\ge1$, including repeated eigenvalues,
\begin{equation}
 \delta_N(s)=N^{-2}\sum_{m=1}^{N-1}\widehat H_s(m/N)|T_m|^2
   +\frac1{2N}\sum_{j\ne k}\sum_{\ell\in\mathbb Z}
                     q_s(x_j-x_k+\ell N).
 \label{eq:abel-positive-20261016}
\end{equation}
Both terms are nonnegative. Equality holds precisely for a rotated clock.
\end{theorem}
\begin{proof}
The transform of $q_s$ is $(s^2/\sinh^2s)$ times the convolution of
$(1-|t|)_+$ and $e^{-2s|t|}$. Elementary integration gives the displayed
transform of $H_s$. Periodize all kernels at the integer period $N$.
The Fourier series are absolutely convergent, and the $H_s$ series has
finite support. The diagonal of the $q_s$ sum contributes exactly
$q_s(0)/2$, since $q_s(\ell N)=0$ for nonzero integers $\ell$.
Finally $(\widehat H_s(0)+q_s(0)-1)/2=(e^{2s}-1)^{-1}$.
\end{proof}

\subsection{An optimal square-root stability theorem}
Let $d_R$ be the bounded local semimetric obtained by taking the supremum
of $|\int\phi\,d(\mu-\nu)|$ over functions supported on $[-R,R]$ with
supremum norm and Lipschitz constant at most one, and then truncating at one.
\begin{theorem}[Stationary clock rigidity]
Let $\Xi$ be a stationary counting process on $\mathbb R$ with intensity
one, allowing multiplicities and an empty component. Suppose its ordered
$q_s$ pair-energy density is $e<\infty$. For every fixed $R$,
\[
 W_{d_R}\bigl(\mathcal L(\Xi),\mathcal L(\mathbb Z+V)\bigr)
       \le \min\{1,C_{s,R}\sqrt e\},
 \qquad V\sim\operatorname{Unif}[0,1).
\]
For the uniform-origin unfolded spectrum, $e\le2\delta_N(s)$, so
$W_{d_R}\le\min\{1,C'_{s,R}\sqrt{\delta_N(s)}\}$ uniformly in $N$.
The exponent $1/2$ is optimal already at $R=2$ over these spectra.
\end{theorem}
\begin{proof}
Let $\nu$ be the intensity-normalized outgoing-gap law, retaining zero
gaps between coincident particles. Coverage of the line by the positive
gaps gives $\int g\,d\nu=1-p_0$, where $p_0$ is the empty-component
probability. Consequently
\[
 A:=\int(1-g)_+\,d\nu
   =\Pr\{\Xi([0,1))=0\},\qquad
 \int|g-1|\,d\nu\le2A.
\]
The chord bound $q_s(g)\ge c_s(1-g)^2$ on $[0,1]$, where
$c_s=4s^3/[\pi^2\sinh^2s(s^2+\pi^2)]$, yields $A\le\sqrt{e/c_s}$.
On a fixed enlarged window, the expected sum of incoming and outgoing
gap deviations is $O_R(A)$. Unless that sum exceeds $1/4$ or the anchor
unit interval is empty, consecutive-index matching with the anchor clock
has total local displacement $O_R(A)$ in expectation. The boundary gaps
are included in the marks. Translating the matched configurations by an
independent uniform variable preserves the stationary process law and
makes the clock phase exactly uniform. This proves the estimate without
a separate count-integrability premise.

For optimality take even $N$ and $x_j=j+\epsilon(-1)^j$. Writing $z=e^{-s}$,
\[
 \delta_N(s)=\frac{z(1-z)(1-\cos(2\pi\epsilon))}
 {2(1+z)(1+2z\cos(2\pi\epsilon)+z^2)}\asymp_s\epsilon^2.
\]
Every translate of this configuration is at $d_2$ distance at least
$\epsilon$ from every clock: two consecutive points in $[-1,1)$ have
fractional separation $2\epsilon$, detected by the nonnegative admissible
test $\min\{\operatorname{dist}(x,\mathbb Z+t),(2-|x|)_+,1\}$.
\end{proof}

The CUE excess is positive. The limiting mean low-mode part of
\eqref{eq:abel-positive-20261016} is
\[
 I_s=\frac{\sinh(2s)-2s-4s^3/3}{8s^2\sinh^2s},
\]
and the geometric part tends to
$\tfrac12\int q_s(x)[1-(\sin(\pi x)/(\pi x))^2],dx$.
Neither limit is established here for the physical fixed map. Vanishing
excess would instead force its clock law.

\subsection{All sublinear powers and the Heisenberg exponential rate}
Let $c(\nu)$ count rooted cyclic words with directed-edge table $\nu$, and put
\[
 B_{N,m}=N^{-m}\sum_\nu c(\nu)^2,\qquad
 G_{N,m}=N^{-m}\sum_\nu c(\nu)^2\prod_{a,b}\nu_{ab}!.
\]
For an exactly flat unitary, $B_{N,m}$ is its positive equal-edge sector.
The classical Gaussian benchmark is
\[
 G_{N,m}=\frac{\prod_{j=0}^m(N+j)-\prod_{j=0}^m(N-j)}
                     {(m+1)N^m},
\]
with falling product zero for $m\ge N$ \cite{KPSS2002}.
\begin{theorem}[Flat equal-edge comparison and collision mean]
For $N\ge3$ and $3\le m\le N$,
\[
 0\le1-B_{N,m}/G_{N,m}
       \le\frac{m(m-1)}{N(N+1)}+\frac{15m}{N^2}.
\]
At $m=1$ the deficit is zero and at $m=2$ it is $1/(2N)$. Thus
$B_{N,m}/G_{N,m}\to1$ for every $m=o(N)$. For $m/N\to\alpha>0$,
\[
 \liminf B_{N,m}/G_{N,m}\ge2^{-2\alpha^2},\qquad
 \frac{\log B_{N,m}}N\longrightarrow
                 (1+\alpha)\log(1+\alpha)-\alpha.
\]
Under the exact Wick table law proportional to $c(\nu)^2\prod\nu_{ab}!$,
\[
 \mathbb E_W\sum_{a,b}\binom{\nu_{ab}}2\longrightarrow2\alpha^2,
 \qquad \Pr_W\{\max_{a,b}\nu_{ab}\ge3\}\longrightarrow0.
\]
For $m\ge N+2$ the finite mean is exactly
\[
 F(N,m)=\frac{2m\{(N+5)m(m+1)-(N+2)(N+3)\}}
                    {(N+1)(N+2)(N+3)(m+2)}.
\]
On every fixed positive compact range of $m/N$, the error from this rational
expression is exponentially small in $N$. With $\alpha_N=m/N$ the mean is
$2\alpha_N^2-2(\alpha_N^2+\alpha_N+1)/N+O(N^{-2})$.
\end{theorem}
\begin{proof}
For unscaled Ginibre $Y$, put $p_m=\operatorname{Tr}Y^m$ and
$h_m=\mathbb E|p_m|^2$. Gaussian Fock adjointness makes the repeated-edge
mean a self-coordinate Hessian norm divided by $2h_m$; the total Hessian
norm is $m(m-1)h_m$. Two sphere fourth moments control off-diagonal
coordinates. The 24 diagonal eighth-moment pairings reduce to four
explicit contractions. The essential contracted matrix satisfies
\[
 \|S\|_{L^2}^2=\langle p_m,\operatorname{Tr}(Y^2D^2)p_m\rangle
       \le m(N+m-1)h_m.
\]
Here $D_{ij}=\partial/\partial Y_{ji}$. The bound follows from hook Schur
orthogonality and the finite identity
$\sum_{a=1}^{m-1}p_ap_{m-a}
=\sum_j(-1)^j(m-1-2j)s_{(m-j,1^j)}$.
The exact partial upper-index swap of the Hessian gives the collision
mean limit. Sixth sphere moments make triple edges negligible.
Keeping the exact 24 and six pairing sums yields the finite hook formula;
stars-and-bars sums of the untruncated hook weights give $F(N,m)$.
The discarded hook tail is exponentially small at $m/N$ in a positive
compact interval. Jensen then bounds $B/G$ away from zero, transferring
the classical Gaussian exponential rate.
\end{proof}

The following continuation, \texttt{BGS-TRACE-094}, closes the limiting
collision distribution and exact flat ratio. Its high-derivative argument
is an additional input beyond the first-moment calculation above.
For a flat unitary, unitarity forces the remaining signed sector to equal
$-B_{N,m}+O(N^2)$. Obtaining its $O(N)$ CUE remainder, extending the sharp
flat comparison to the actual non-flat profile, and proving the full
every-order signed hierarchy remain distinct unresolved obligations.

\subsection{Complete Wick collision law at Heisenberg powers}
\label{subsec:wick-poisson-20261016}
The canonical proof is
\path{mathematics/wick_collision_poisson_20261016.md}.
Let $\mathbb W(\nu)=c(\nu)^2\prod_e\nu_e!/h_m$ and put
$C_N=\sum_e\binom{\nu_e}{2}$ and $D_N=\#\{e:\nu_e\ge2\}$.
Directed loops are included.
\begin{theorem}[Complete flat collision comparison]
For every fixed $0<a<b<\infty$ and positive integer $r$, uniformly for
$aN\le m\le bN$,
\[
 \mathbb E_{\mathbb W}(D_N)_r
    =[2(m/N)^2]^r+O_{a,b,r}(N^{-1}),\qquad
 \Pr_{\mathbb W}\{\max_e\nu_e\ge3\}=O_{a,b}(N^{-1}).
\]
Consequently, if $m/N\to\alpha>0$,
\[
 C_N\ \Longrightarrow\ \operatorname{Poisson}(2\alpha^2),
 \qquad \frac{B_{N,m}}{G_{N,m}}\longrightarrow e^{-\alpha^2}.
\]
This is a collision law of an auxiliary positive measure, not an
energy-level spacing law.
\end{theorem}
\begin{proof}
For the $k$th derivative tensor
$H^{(k)}=(\partial_{a_1b_1}\cdots\partial_{a_kb_k}p_m)$, every partial matching
$P$ of upper and lower indices obeys
\begin{equation}
 \|C_PH^{(k)}\|_2^2\le C_km(N+m)^{k-1}h_m
          \le C_{a,b,k}(m)_kh_m.
 \label{eq:wick-contraction-20261016}
\end{equation}
Here $k$ is fixed and $N$ is sufficiently large. To prove the first bound,
expand the marked trace into $(k-1)!$ cyclic orders and compositions
$r_1+\cdots+r_k=m-k$. Move all $k$ derivatives in one factor of the
squared norm to creation variables by Fock adjointness. The resulting
index network consists of $k$ unit creation edges and $k$ segments of
lengths $r_i$. Partial matches and the remaining cross-copy links form
directed cycles. Each zero-length cycle consumes a distinct zero $r_i$.
If exactly $q$ gaps are zero there are
$\binom{k}{q}\binom{m-k-1}{k-q-1}$ compositions and at most $q$ factors
$N$ from zero traces. A product of $\ell\le2k$ positive traces has overlap
with $p_m$ at most $2^\ell h_m$: its hook coefficients are coefficients of
$\prod_v(1-t^v)/(1-t)$ and hence have modulus at most $2^\ell$.
Summing over $q\le k-1$ proves \eqref{eq:wick-contraction-20261016}.

Conjugation invariance now permits the classical finite Haar formula
\cite{CollinsSniady2006}. At degree $2k$ its Gram matrix is
$G_{\sigma\tau}=N^{\#\operatorname{cycles}(\sigma^{-1}\tau)}$.
For fixed $k$, $G^{-1}=N^{-2k}(I+O_k(N^{-1}))$ entrywise.
Every input pairing is the inner product of two partial contractions
in \eqref{eq:wick-contraction-20261016}. Therefore, uniformly over all
coordinate choices, including coincidences,
\[
 \frac{\|\partial_{e_1}\cdots\partial_{e_k}p_m\|_2^2}{h_m}
       \le C_{a,b,k}N^{-k}.
\]
If the prescribed upper and lower label sets are disjoint, the leading
pairings are fully cross. Upper multiplicities $A_u$ and lower
multiplicities $B_v$ give $\prod_u A_u!\prod_v B_v!$ such pairings.
For $k\ge2$, mixed Fock homogeneity transfers the Hessian swap identity exactly:
\[
 K_k=(m-2)_{k-2}K_2,\qquad
 K_k/\|H^{(k)}\|_2^2=1-2\bar j/(m-1),\qquad
 \bar j\le(N-1)/m.
\]
A bounded number of upper transpositions makes each leading overlap
$1+O_{a,b,k}(N^{-1})$ times the total energy. Thus the coordinate norm
equals $(m)_kN^{-2k}[\prod A_u!\prod B_v!+O_{a,b,k}(N^{-1})]h_m$.

The exact coefficient identity identifies these derivative norms with
$\mathbb E_{\mathbb W}\prod_e(\nu_e)_{k_e}$.
For $r$ distinct edges each used twice, with all $2r$ endpoints distinct,
the multiplicity is $4^r$ and the number of tuples is $(N)_{2r}$.
Dividing by $2^r$ gives the asserted leading factorial moment. All
other tuples number $O_r(N^{2r-1})$ and the uniform coordinate bound
makes their total $O(N^{-1})$. Replacing binomial counts by indicators
costs a sum of derivative order $2r+1$, again $O(N^{-1})$.
Order three bounds the probability of any triple edge by $O(N^{-1})$.
Finite Bonferroni bounds, with truncation removed only after the limit
in $N$, give $\mathbb E(1-q)^{D_N}\to e^{-2\alpha^2q}$ for $0\le q\le1$.
Finally $C_N=D_N$ and $\prod\nu_e!=2^{C_N}$ outside the triple-edge
event. These bounded transforms prove both conclusions.
\end{proof}

The same argument gives independent Poisson counts in arbitrary disjoint
deterministic edge classes of limiting densities $\beta_j$, with means
$2\alpha^2\beta_j$. For fixed distinct vertices, the outgoing counts
$n_a=\sum_b\nu_{ab}$ instead have independent geometric limits of mean
$\alpha$: their mixed factorial moments tend to
$\prod_i k_i!\alpha^{\sum_i k_i}$. Sum the coordinate derivative formula
over the lower labels; all distinct lower labels outside the fixed upper
set give this main term, and the remaining $O(N^{K-1})$ choices are
controlled uniformly. Exchangeability also gives convergence in probability
of the empirical vertex-occupation distribution. Incoming and outgoing
counts of one vertex remain equal exactly.

As a separate corollary, let $Z_{ab}=\xi_{ab}/\sqrt N$ be iid circular
complex entries with $\mathbb E|\xi|^2=1$ and
$\mathbb E|\xi|^{2j}\le j!$ for every $j$. With
$\kappa=\mathbb E|\xi|^4$,
\[
 \frac{\mathbb E|\operatorname{Tr}Z^m|^2}{G_{N,m}}
       \longrightarrow e^{(\kappa-2)\alpha^2}.
\]
Indeed the exact ratio is $\mathbb E_{\mathbb W}\prod_e q_{\nu_e}$,
$q_j=\mathbb E|\xi|^{2j}/j!\in[0,1]$. Away from triple edges this is
$(\kappa/2)^{C_N}$. The bounded transform limit applies. This proof
does not cover $\kappa>2$ or the non-flat deterministic physical profile.
